% ================== 2. Related concepts / background ==================
\section{Related Concepts and Theoretical Background}\label{sec:background}

\subsection{Temporal memory and state sufficiency}\label{sec:bkg-state}

A dynamical system is Markovian in a state variable if that variable
summarizes all information from the past that is relevant to the future.
For leaky integrate-and-fire (LIF) neurons
\citep{lapicque1907,abbott1999,burkitt2006}, driven by an external
scalar input, the membrane potential together with the input history
constitutes the state; with a fixed input schedule the membrane potential
alone is a sufficient statistic in the sense that identical potentials
with identical future inputs produce identical futures.  Working-memory
models at the population level instead maintain content through
attractor or delay activity \citep{goldmanrakic1995,durstewitz2000,
compte2000,hopfield1982}, while recurrent sequence models exhibit
temporal state through hidden dynamics \citep{elman1990,maass2002,
jaeger2001}.  For a single neuron, whether history enters the transition
law at all, and whether the scalar potential remains a sufficient
projection, is an empirical question that a state-collision experiment
can answer directly: two histories that lead to (nearly) identical
membrane potentials are injected with an identical probe; if the
next-state potentials diverge, the scalar potential is not a sufficient
state and the projection onto it is not a consistent (lumpable)
reduction \citep{kemeny1960}.

\subsection{Attention as input-dependent weighting}\label{sec:bkg-attn}

Attention in sequence models computes a normalized, input-dependent
weighting of values by query--key compatibility
\citep{vaswani2017,bahdanau2015}; the weights form a Gibbs distribution,
connecting such weighting to statistical mechanics
\citep{jaynes1957,gibbs1902}.  In sensory neuroscience, saliency
weighting describes bottom-up prominence of inputs
\citep{koch1985,itti1998}, and normalization models describe
input-dependent gain \citep{reynolds2009}.  Neuromodulatory gain control
offers a biological analogue of input-conditioned modulation
\citep{astonjones2005,yudayan2005}.  Predictive-coding accounts treat
deviation-from-expectation signals as the driver of cortical gain
\citep{rao1999,srinivasan1982,friston2005,keller2018,bastos2012}, and
single-cell asymmetric response properties are natural substrates for
such deviations \citep{laughlin1981}.  TAN instantiates input-dependent
weighting as a surprise-scaled, Boltzmann-style reweighting of its own
temporal window; the key theoretical property we later exploit is that
the surprise variable is a scalar that multiplies every window entry
uniformly.

\subsection{Routing and content-addressable selection}\label{sec:bkg-routing}

Routing is the ability to direct information from a selected source to a
consumer based on control signals.  Classical proposals in vision route
attended locations into an invariant representation
\citep{olshausen1993}, mixture-of-experts models route inputs to experts
through a gating network \citep{jacobs1991,shazeer2017,fedus2022}, and
mechanistic analyses of transformers interpret attention heads as
content-addressable key--value lookups that implement induction and
copying behaviour \citep{elhage2021,olsson2022}.  Binding -- linking a
selected content to a role or query -- has been associated with
feature-integration \citep{treisman1980}, correlation/synchrony
\citep{von_der_malsburg1995,von_der_malsburg1986,singer1995,fries2005},
and, in cellular terms, with dendritic and circuit mechanisms
\citep{london2005,larkum2013,spruston2008}.  An important interpretive
caution applies to all of these literatures: observed weight values are
not by themselves demonstrations of routing or explanation
\citep{jainwallace2019,serranosmith2019,wiegreffe2019}; a weighting
mechanism must be shown to change \emph{which content} is selected as a
function of a control signal, not merely how strongly a fixed content is
weighted.

\subsection{Mechanistic identifiability}\label{sec:bkg-ident}

When a task can be solved through unintended shortcuts -- fixed
positional memory, amplitude cues, or query--label correlations -- a
measured advantage of one architecture over another cannot be attributed
to the mechanism under study.  Mechanistic identifiability requires that
(i) simpler storage/integration mechanisms cannot solve the task through
linear readouts, (ii) the mechanism of interest is observable (e.g.\ the
attention weights themselves), and (iii) counterfactual manipulations
(permuting positions, values, queries) change the mechanism in the
direction predicted by the scientific claim.  This paper applies this
audit discipline to every layer of the hierarchy above; negative audit
outcomes are treated as mechanistic results, not as experimental
failures.
